\Subhead{6.3.2}{Реализуемые линейные фильтры и уравнение Винера--Хопфа}

Реализуемый без предвидения фильтр использует только $u\geq0$:
\begin{equation}\label{eq:wienerhopf}
 \int_0^\infty h(u)C_y(\tau-u)\dd u=C_{xy}(\tau),\qquad \tau\geq0.
\end{equation}
Это уравнение Винера–Хопфа. Ортогональность требуется только к доступному прошлому. Поэтому нельзя просто обнулить отрицательную часть импульсной характеристики некаузального оптимума: оставшаяся часть в общем случае тоже должна измениться.

Смысл решения через спектральную факторизацию: представить $S_y=F F^*$, где $F$ — устойчивый причинный минимально-фазовый формирующий фильтр с допустимым обратным. Деление наблюдения на $F$ выбеливает его. Для белого входа оптимальные веса равны корреляции с оцениваемым сообщением на доступных задержках; затем возвращаются к исходному наблюдению. Формально
\begin{equation}
 H_{\rm c}(\omega)=\frac1{F(\omega)}
 \left[\frac{S_{xy}(\omega)}{F^*(\omega)}\right]_+,
\end{equation}
где $[\cdot]_+$ оставляет причинную часть обратного преобразования. Нужны условия существования факторизации и конечности ошибки. В дискретном конечномерном приближении задача становится линейной системой с ковариационной матрицей.

Для прогноза меняется взаимная корреляция желаемого будущего $x(t+T)$ с прошлым $y$; для сглаживания разрешаются также будущие данные. Причинное ограничение не является дополнительным шумом: оно ограничивает доступную информацию.
